\documentclass[11pt]{article}

\usepackage[a4paper,margin=1in]{geometry}
\usepackage{amsmath,amssymb,amsfonts,amsthm,mathtools}
\usepackage{enumitem}
\usepackage{hyperref}
\usepackage{mathrsfs}

\hypersetup{
  colorlinks=true,
  linkcolor=blue,
  citecolor=blue,
  urlcolor=blue
}

\newtheorem{theorem}{Theorem}[section]
\newtheorem{proposition}[theorem]{Proposition}
\newtheorem{lemma}[theorem]{Lemma}
\newtheorem{corollary}[theorem]{Corollary}
\newtheorem{definition}[theorem]{Definition}
\newtheorem{remark}[theorem]{Remark}

\newcommand{\ZZ}{\mathbb Z}
\newcommand{\NN}{\mathbb N}
\newcommand{\BB}{\mathbb B}
\newcommand{\QQ}{\mathbb Q}
\newcommand{\CC}{\mathbb C}
\newcommand{\Spec}{\operatorname{Spec}}
\newcommand{\Hom}{\operatorname{Hom}}
\newcommand{\Diag}{\operatorname{Diag}}
\newcommand{\Id}{\operatorname{Id}}
\newcommand{\Frac}{\operatorname{Frac}}
\newcommand{\Cl}{\operatorname{Cl}}
\newcommand{\supp}{\operatorname{supp}}
\newcommand{\coker}{\operatorname{coker}}
\newcommand{\im}{\operatorname{im}}
\newcommand{\Ker}{\operatorname{Ker}}
\newcommand{\ord}{\operatorname{ord}}
\newcommand{\tr}{\operatorname{tr}}

\title{
A Split-Zero Repair Dossier:\\
Support-Complete Proofs for \(G(R)=R\sqcup\{\tau\}\),\\
Boundary Prime Structure, Semimodules, Zeta Sheets,\\
and Residual Divisor Certificates
}

\author{The Clankers / prepared in continuation}
\date{May 2026}

\begin{document}
\maketitle

\begin{abstract}
We give a support-complete writeout of the split-zero proofs that must be used whenever an argument is carried out over
\[
G(R)=R\sqcup\{\tau\}.
\]
The semiring zero is \(\tau\), while the supported arithmetic zero is \(e=0_R\in R\subset G(R)\). This distinction is not cosmetic: when \(R\) is a domain, \(e\) is reducible and divisibility-prime. Consequently, classical factorization arguments are conservative only on the nonzero supported sheet \(R\setminus\{0_R\}\), while any argument involving zero, coprimality, primehood, zeta factors, support, semimodules, or residual divisors must be retyped over the split support structure. We prove the normal form, universal quotient, support character, ideals, prime ideals, congruences, localizations, conservative ideal arithmetic, non-Eulerian separated zeta sheet, semimodule classification, cancellative reduction, Boolean free semimodules, projective support diagrams, shifted zeta deformation, boundary coprimality matrices, and the corrected split-zero formulation of residual Erdős--Straus divisor certificates.
\end{abstract}

\tableofcontents

\section{Standing convention}

Throughout \(R\) denotes a nonzero commutative unital ring.

\begin{definition}[Split-zero semiring]
Define
\[
G(R):=R\sqcup\{\tau\},
\]
where \(\tau\notin R\), with operations
\[
r\oplus s=r+s,\qquad r\oplus\tau=\tau\oplus r=r,\qquad \tau\oplus\tau=\tau,
\]
and
\[
rs=rs,\qquad r\tau=\tau r=\tau,\qquad \tau^2=\tau
\]
for \(r,s\in R\).

We write
\[
S:=G(R),\qquad e:=0_R\in R\subset S.
\]
Thus
\[
0_S=\tau,\qquad e=0_R\neq\tau.
\]
\end{definition}

The entire purpose of the split-zero construction is to keep separate the two zero-like objects
\[
\tau
\qquad\text{and}\qquad
e=0_R.
\]
Accordingly, any proof over \(G(R)\) that says merely ``\(x=0\)'' is under-specified. It must say either
\[
x=\tau,
\]
or
\[
x=e=0_R,
\]
or
\[
p_R(x)=0_R
\]
for the ring reflection \(p_R\) defined below.

\section{Normal form and universal maps}

\begin{definition}[Support/value model]
Let
\[
D_R:=\{(r,b)\in R\times\BB:\ b=0\Rightarrow r=0_R\}.
\]
Equivalently,
\[
D_R=\{(0_R,0)\}\cup (R\times\{1\}).
\]
\end{definition}

\begin{theorem}[Algebraic normal form]
There is a canonical semiring isomorphism
\[
\iota:G(R)\xrightarrow{\sim}D_R
\]
defined by
\[
\iota(\tau)=(0_R,0),\qquad \iota(r)=(r,1)\quad(r\in R).
\]
The missing states are precisely
\[
(r,0),\qquad r\neq0_R.
\]
\end{theorem}

\begin{proof}
The map \(\iota\) is plainly bijective.

The subset \(D_R\subset R\times\BB\) is closed under componentwise addition and multiplication. Indeed:
\[
(r,1)+(s,1)=(r+s,1),
\]
\[
(r,1)+(0_R,0)=(r,1),
\]
\[
(0_R,0)+(0_R,0)=(0_R,0),
\]
and
\[
(r,1)(s,1)=(rs,1),
\]
\[
(r,1)(0_R,0)=(0_R,0),
\]
\[
(0_R,0)^2=(0_R,0).
\]
Transporting these operations along \(\iota\) gives exactly the defining operations of \(G(R)\). Hence \(\iota\) is a semiring isomorphism.
\end{proof}

\begin{definition}[Ring reflection and support character]
Define
\[
p_R:G(R)\to R
\]
by
\[
p_R(\tau)=0_R,\qquad p_R(r)=r\quad(r\in R).
\]
Define
\[
\chi_R:G(R)\to\BB
\]
by
\[
\chi_R(\tau)=0,\qquad \chi_R(r)=1\quad(r\in R).
\]
\end{definition}

\begin{corollary}[Two zero fibers]
One has
\[
p_R^{-1}(0_R)=\{\tau,e\},
\]
whereas
\[
\chi_R^{-1}(0)=\{\tau\},
\qquad
\chi_R^{-1}(1)=R.
\]
\end{corollary}

\begin{proof}
By definition,
\[
p_R(\tau)=0_R,\qquad p_R(e)=p_R(0_R)=0_R.
\]
If \(r\in R\) and \(p_R(r)=0_R\), then \(r=0_R=e\). This proves
\[
p_R^{-1}(0_R)=\{\tau,e\}.
\]

Likewise \(\chi_R(\tau)=0\), while \(\chi_R(r)=1\) for every supported element \(r\in R\), including \(r=e=0_R\). Hence
\[
\chi_R^{-1}(0)=\{\tau\}.
\]
\end{proof}

\begin{theorem}[Universal ring quotient]
For every commutative ring \(A\), composition with \(p_R\) induces a bijection
\[
\Hom_{\mathrm{Ring}}(R,A)
\xrightarrow{\sim}
\Hom_{\mathrm{SemiRing}}(G(R),A).
\]
Equivalently, \(R\) is the universal ring quotient of \(G(R)\).
\end{theorem}

\begin{proof}
Let
\[
f:G(R)\to A
\]
be a unital semiring homomorphism, where \(A\) is a ring.

Since \(\tau=0_S\), one has
\[
f(\tau)=0_A.
\]
Since \(e\oplus e=e\), applying \(f\) gives
\[
f(e)+f(e)=f(e).
\]
The additive group of \(A\) is cancellative, so
\[
f(e)=0_A.
\]

Restrict \(f\) to \(R\subset G(R)\). For \(r,s\in R\),
\[
f(r+s)=f(r\oplus s)=f(r)+f(s),
\]
and
\[
f(rs)=f(r)f(s).
\]
Also \(f(1_R)=1_A\). Hence \(f|_R:R\to A\) is a ring homomorphism. Since \(f(\tau)=0_A=f|_R(0_R)\), one has
\[
f=f|_R\circ p_R.
\]

Conversely, if \(g:R\to A\) is a ring homomorphism, define
\[
\widetilde g:G(R)\to A
\]
by
\[
\widetilde g(\tau)=0_A,\qquad \widetilde g(r)=g(r).
\]
The defining addition and multiplication laws of \(G(R)\) are immediately preserved. Hence \(\widetilde g\) is a semiring homomorphism and
\[
\widetilde g=g\circ p_R.
\]
The two constructions are inverse.
\end{proof}

\begin{theorem}[Unique Boolean support character]
The support character
\[
\chi_R:G(R)\to\BB
\]
is the unique unital semiring homomorphism from \(G(R)\) to the Boolean semiring \(\BB\).
\end{theorem}

\begin{proof}
Let
\[
\psi:G(R)\to\BB
\]
be a unital semiring homomorphism.

Since \(\tau=0_S\), one has
\[
\psi(\tau)=0.
\]
Since \(\psi\) is unital,
\[
\psi(1_R)=1.
\]

Inside \(R\), one has
\[
1_R+(-1_R)=0_R=e.
\]
Therefore in \(G(R)\),
\[
1_R\oplus (-1_R)=e.
\]
Applying \(\psi\),
\[
\psi(e)=\psi(1_R)\vee\psi(-1_R)=1\vee\psi(-1_R)=1.
\]

Now take arbitrary \(r\in R\). Since
\[
r=e\oplus r,
\]
we get
\[
\psi(r)=\psi(e)\vee\psi(r)=1\vee\psi(r)=1.
\]
Thus
\[
\psi(\tau)=0,\qquad \psi(r)=1\quad(r\in R),
\]
so
\[
\psi=\chi_R.
\]
\end{proof}

\begin{proposition}[No section from \(\BB\)]
There is no unital semiring homomorphism
\[
\BB\to G(R)
\]
when \(R\) is nonzero.
\end{proposition}

\begin{proof}
Suppose
\[
f:\BB\to G(R)
\]
is unital. Then
\[
f(1)=1_R.
\]
But in \(\BB\),
\[
1+1=1.
\]
Therefore
\[
1_R=f(1)=f(1+1)=f(1)\oplus f(1)=1_R\oplus1_R=1_R+1_R=2\cdot1_R.
\]
Thus
\[
1_R=2\cdot1_R,
\]
so
\[
1_R=0_R,
\]
contradicting that \(R\) is nonzero.
\end{proof}

\section{Ideals, prime ideals, and dimension}

\begin{theorem}[Ideal classification]
The ideals of \(G(R)\) are exactly
\[
\{\tau\}
\]
and
\[
I_J:=J\cup\{\tau\},
\]
where \(J\triangleleft R\) is an ordinary ideal.

For \(R=\ZZ\), every ideal is either
\[
\{\tau\}
\]
or
\[
I_n:=n\ZZ\cup\{\tau\}
\qquad(n\ge0).
\]
\end{theorem}

\begin{proof}
Let \(I\subseteq S=G(R)\) be an ideal. Since \(\tau=0_S\), one has
\[
\tau\in I.
\]
Set
\[
J:=I\cap R.
\]

If \(J=\varnothing\), then \(I=\{\tau\}\).

Assume \(J\neq\varnothing\). We prove \(J\triangleleft R\). Let \(a,b\in J\). Then \(a,b\in I\), and since \(I\) is closed under \(\oplus\),
\[
a\oplus b=a+b\in I.
\]
Because \(a+b\in R\), it follows that
\[
a+b\in J.
\]

Let \(a\in J\). Since \(-1_R\in R\subset S\), multiplicative absorption gives
\[
(-1_R)a=-a\in I.
\]
Again \(-a\in R\), hence
\[
-a\in J.
\]
Thus \(J\) is an additive subgroup of \(R\).

If \(r\in R\) and \(a\in J\), then
\[
ra\in I
\]
by absorption. Since \(ra\in R\), one has
\[
ra\in J.
\]
Therefore \(J\triangleleft R\).

Moreover,
\[
I=J\cup\{\tau\}.
\]

Conversely, if \(J\triangleleft R\), then \(J\cup\{\tau\}\) is closed under addition:
two supported elements add inside \(J\), while adding \(\tau\) returns the other summand. It is multiplicatively absorbing: products with \(\tau\) are \(\tau\), and products of supported elements with elements of \(J\) remain in \(J\). Therefore \(J\cup\{\tau\}\) is an ideal.

For \(R=\ZZ\), every ideal has the form \(n\ZZ\), \(n\ge0\), giving the stated classification.
\end{proof}

\begin{theorem}[Prime ideal classification]
The prime ideals of \(G(R)\) are:
\[
P_\tau:=\{\tau\},
\]
and
\[
P_{\mathfrak p}:=\mathfrak p\cup\{\tau\},
\]
where \(\mathfrak p\) is a prime ideal of \(R\).

For \(R=\ZZ\), the prime ideals are precisely
\[
P_\tau=\{\tau\},
\]
\[
P_0=\{0,\tau\},
\]
and
\[
P_p=p\ZZ\cup\{\tau\}
\]
for ordinary rational primes \(p\).
\end{theorem}

\begin{proof}
First \(P_\tau=\{\tau\}\) is prime. If
\[
xy\in P_\tau,
\]
then
\[
xy=\tau.
\]
If both \(x\) and \(y\) were supported, then \(x,y\in R\) and their product would lie in \(R\), not in \(\{\tau\}\). Hence at least one of \(x,y\) equals \(\tau\). Thus one of \(x,y\) lies in \(P_\tau\).

Now let \(J\triangleleft R\) and \(I_J=J\cup\{\tau\}\).

Assume \(I_J\) is prime. If \(a,b\in R\) and \(ab\in J\), then \(ab\in I_J\). Hence
\[
a\in I_J
\quad\text{or}\quad
b\in I_J.
\]
Since \(a,b\in R\), this means
\[
a\in J
\quad\text{or}\quad
b\in J.
\]
Therefore \(J\) is prime in \(R\).

Conversely, suppose \(J=\mathfrak p\) is prime in \(R\). Let \(x,y\in G(R)\) and suppose
\[
xy\in\mathfrak p\cup\{\tau\}.
\]
If \(x=\tau\) or \(y=\tau\), then one factor already lies in \(P_{\mathfrak p}\). Otherwise \(x,y\in R\). Then \(xy\in R\), so
\[
xy\in\mathfrak p.
\]
Since \(\mathfrak p\) is prime,
\[
x\in\mathfrak p
\quad\text{or}\quad
y\in\mathfrak p.
\]
Thus
\[
x\in P_{\mathfrak p}
\quad\text{or}\quad
y\in P_{\mathfrak p}.
\]
Therefore \(P_{\mathfrak p}\) is prime.

For \(R=\ZZ\), the prime ideals of \(\ZZ\) are \((0)\) and \((p)\) for ordinary primes \(p\). This gives the listed prime ideals.
\end{proof}

\begin{corollary}[Dimension of \(G(\ZZ)\)]
\[
\dim G(\ZZ)=2.
\]
More precisely, every maximal prime chain has the form
\[
P_\tau\subsetneq P_0\subsetneq P_p.
\]
\end{corollary}

\begin{proof}
The inclusions
\[
\{\tau\}\subsetneq\{0,\tau\}\subsetneq p\ZZ\cup\{\tau\}
\]
are immediate.

By the prime ideal classification, there are no further prime ideals. Distinct \(P_p\) and \(P_q\) are incomparable when \(p\neq q\). Therefore no chain has length greater than \(2\).
\end{proof}

\section{Divisibility, units, irreducibles, and the supported-zero prime}

\begin{theorem}[Units]
If \(R\) is an integral domain, then
\[
G(R)^\times=R^\times.
\]
In particular,
\[
G(\ZZ)^\times=\{\pm1\}.
\]
\end{theorem}

\begin{proof}
Let \(u\in G(R)\) be a unit. Then there exists \(v\in G(R)\) with
\[
uv=1.
\]
If \(u=\tau\), then
\[
uv=\tau,
\]
contradiction. Thus \(u\in R\). Similarly \(v\in R\). Therefore \(uv=1\) in \(R\), so \(u\in R^\times\).

Conversely, every unit of \(R\) is plainly a unit of \(G(R)\).
\end{proof}

\begin{theorem}[Nonzero supported irreducibles are classical]
Assume \(R\) is an integral domain. Let
\[
x\in R\setminus\{0_R\}.
\]
Then \(x\) is irreducible in \(G(R)\) if and only if \(x\) is irreducible in \(R\).
\end{theorem}

\begin{proof}
Suppose
\[
x=ab
\]
in \(G(R)\). Since \(x\in R\setminus\{0_R\}\), neither \(a\) nor \(b\) can be \(\tau\), because any product involving \(\tau\) is \(\tau\). Therefore \(a,b\in R\), and the factorization is an ordinary factorization in \(R\).

By the preceding theorem, the units of \(G(R)\) are the units of \(R\). Hence the irreducibility condition is exactly the same in \(G(R)\) and in \(R\).
\end{proof}

\begin{corollary}[Irreducibles of \(G(\ZZ)\)]
The irreducible elements of \(G(\ZZ)\) are exactly
\[
\pm p,
\]
where \(p\) ranges over ordinary rational primes.
\end{corollary}

\begin{proof}
Apply the previous theorem to \(R=\ZZ\), whose nonzero irreducibles are exactly \(\pm p\).
\end{proof}

\begin{theorem}[Supported zero is reducible and divisibility-prime]
Assume \(R\) is an integral domain. Then
\[
e=0_R\in G(R)
\]
is reducible and prime in the divisibility sense:
\[
e\mid xy
\quad\Longrightarrow\quad
e\mid x\ \text{or}\ e\mid y.
\]
Consequently, prime does not imply irreducible in \(G(R)\).
\end{theorem}

\begin{proof}
First \(e\) is not a unit, since \(0_R\notin R^\times\). Moreover,
\[
e=e\cdot e,
\]
so \(e\) is reducible.

Now suppose
\[
e\mid xy.
\]
Then there exists \(c\in G(R)\) such that
\[
xy=ec.
\]
If \(c=\tau\), then
\[
ec=e\tau=\tau.
\]
If \(c\in R\), then
\[
ec=0_Rc=0_R=e.
\]
Therefore
\[
xy\in\{\tau,e\}.
\]

If \(x=\tau\), then
\[
x=\tau=e\tau,
\]
so \(e\mid x\). Similarly if \(y=\tau\), then \(e\mid y\).

Assume now \(x,y\neq\tau\). Then \(x,y\in R\). Since \(xy\in R\) and \(xy\in\{\tau,e\}\), we must have
\[
xy=e=0_R.
\]
Because \(R\) is an integral domain,
\[
x=0_R
\quad\text{or}\quad
y=0_R.
\]
Thus \(x=e\) or \(y=e\). If \(x=e\), then \(x=e\cdot1\), so \(e\mid x\). Likewise if \(y=e\), then \(e\mid y\).

Therefore \(e\) is divisibility-prime.

Since \(e\) is prime in this divisibility sense but reducible, prime does not imply irreducible in \(G(R)\).
\end{proof}

\begin{corollary}[Failure of unique factorization as a semiring]
If \(R\) is an integral domain, then \(G(R)\) is not a unique factorization semiring in any definition requiring prime elements to be irreducible or factorizations to be controlled by irreducibles.
\end{corollary}

\begin{proof}
The element \(e=0_R\) is divisibility-prime but reducible. Hence the usual implication from prime to irreducible fails. Moreover, \(e=e^2=e^3=\cdots\) gives nonterminating supported-zero factorizations. Thus any unique factorization theory that does not explicitly exclude \(e\) fails.
\end{proof}

\section{Congruences and quotient sheets for \(G(\ZZ)\)}

Now let
\[
S=G(\ZZ).
\]

\begin{definition}[Two congruence families]
For \(n\ge0\), define \(\rho_n\) by:
\[
a\rho_n b
\quad\Longleftrightarrow\quad
a\equiv b\pmod n
\qquad(a,b\in\ZZ),
\]
and \(\tau\) is equivalent only to itself.

Define \(\rho_n'\) by:
\[
a\rho_n' b
\quad\Longleftrightarrow\quad
a\equiv b\pmod n
\qquad(a,b\in\ZZ),
\]
and additionally
\[
\tau\rho_n'0.
\]
Equivalently, the \(\rho_n'\)-class of \(\tau\) is
\[
n\ZZ\cup\{\tau\}.
\]
\end{definition}

\begin{theorem}[Congruence classification]
Every congruence on \(S=G(\ZZ)\) is exactly one of
\[
\rho_n
\]
or
\[
\rho_n'
\]
for a unique \(n\ge0\). Moreover,
\[
S/\rho_n\cong G(\ZZ/n\ZZ),
\]
whereas
\[
S/\rho_n'\cong\ZZ/n\ZZ.
\]
For \(n=0\),
\[
S/\rho_0'\cong\ZZ.
\]
\end{theorem}

\begin{proof}
Let \(\sim\) be a congruence on \(S\). Its restriction to \(\ZZ\subset S\) is a ring congruence on \(\ZZ\), so there is a unique \(n\ge0\) such that for \(a,b\in\ZZ\),
\[
a\sim b
\quad\Longleftrightarrow\quad
a\equiv b\pmod n.
\]

There are two cases.

\emph{Case 1: \(\tau\not\sim0\).} We claim \(\tau\) is equivalent to no integer. Suppose \(m\in\ZZ\) and \(m\sim\tau\). Add \(-m\) to both sides:
\[
m\oplus(-m)\sim\tau\oplus(-m).
\]
The left side is \(0\). The right side is \(-m\). Thus
\[
0\sim -m.
\]
Since \(-m\in\ZZ\), this implies \(-m\equiv0\pmod n\), hence \(-m\sim0\). By transitivity,
\[
\tau\sim0,
\]
contradiction. Therefore \(\tau\) is singleton. The remaining classes are the residue classes modulo \(n\), so
\[
\sim=\rho_n.
\]
The quotient has one unsupported zero class and \(n\) supported residue classes, hence is \(G(\ZZ/n\ZZ)\).

\emph{Case 2: \(\tau\sim0\).} Then for \(m\in\ZZ\),
\[
m\sim\tau
\quad\Longleftrightarrow\quad
m\sim0
\quad\Longleftrightarrow\quad
n\mid m.
\]
Hence the class of \(\tau\) is
\[
n\ZZ\cup\{\tau\}.
\]
The quotient is exactly the ordinary ring quotient
\[
\ZZ/n\ZZ.
\]
Thus
\[
\sim=\rho_n'.
\]
\end{proof}

\begin{corollary}[Classical zeta sheet]
Let
\[
\zeta_S(s):=\sum_I N(I)^{-s},
\]
where \(I\) ranges over ideals of finite norm and
\[
N(I):=|S/\rho_I|
\]
for the congruence naturally associated to \(I\). Then
\[
\zeta_S(s)=\zeta(s).
\]
\end{corollary}

\begin{proof}
The ideals of \(S\) are \(\{\tau\}\) and \(I_n=n\ZZ\cup\{\tau\}\).

The ideal \(\{\tau\}\) has quotient \(S\), hence infinite norm. The ideal \(I_0=\{0,\tau\}\) has quotient \(\ZZ\), hence infinite norm.

For \(n\ge1\), the associated quotient is
\[
\ZZ/n\ZZ,
\]
so
\[
N(I_n)=n.
\]
Therefore
\[
\zeta_S(s)=\sum_{n\ge1}n^{-s}=\zeta(s).
\]
\end{proof}

\begin{corollary}[Separated quotient-size sheet]
The support-separated quotient family has finite sizes
\[
|G(\ZZ/n\ZZ)|=n+1.
\]
Hence its Dirichlet series is
\[
\zeta_S^{\mathrm{sep}}(s)
=
\sum_{n\ge1}(n+1)^{-s}
=
\zeta(s)-1.
\]
\end{corollary}

\begin{proof}
The quotient \(G(\ZZ/n\ZZ)\) has \(n\) supported residue classes and one unsupported zero class. Therefore
\[
|G(\ZZ/n\ZZ)|=n+1.
\]
Thus
\[
\sum_{n\ge1}(n+1)^{-s}
=
\sum_{m\ge2}m^{-s}
=
\zeta(s)-1.
\]
\end{proof}

\section{Localization}

\begin{theorem}[Localization away from supported zero]
Let \(M\subset R\) be a multiplicative subset with \(0_R\notin M\). Then
\[
G(R)_M\cong G(R_M).
\]
\end{theorem}

\begin{proof}
Define
\[
\Phi:G(R)_M\to G(R_M)
\]
by
\[
\Phi\left(\frac{\tau}{m}\right)=\tau,
\]
and
\[
\Phi\left(\frac{r}{m}\right)=\frac rm
\qquad(r\in R,\ m\in M).
\]

We prove well-definedness. In semiring localization, two fractions \((x,m)\), \((y,n)\) are equivalent if there exists \(k\in M\) such that
\[
kxn=ky m
\]
in \(G(R)\).

Suppose \(x=\tau\). Then
\[
kxn=k\tau n=\tau.
\]
If \(y\in R\), then \(kym\in R\), never \(\tau\). Hence \(y=\tau\). Therefore all fractions represented by \(\tau\) map consistently to \(\tau\).

If \(x,y\in R\), then the equality is the ordinary ring-localization equality
\[
kxn=kym
\]
in \(R\). This is exactly the equivalence relation defining \(R_M\).

The map \(\Phi\) preserves addition and multiplication by direct calculation on representatives. It is surjective because every element of \(G(R_M)\) is either \(\tau\) or a fraction \(r/m\). It is injective by the same equivalence calculation.

Therefore \(\Phi\) is an isomorphism.
\end{proof}

\begin{theorem}[Total supported localization]
If one localizes \(G(R)\) by inverting every supported element of \(R\), including \(e=0_R\), then the result is the Boolean semiring:
\[
G(R)_R\cong\BB.
\]
\end{theorem}

\begin{proof}
In the localization, \(e=0_R\) is invertible. Hence there exists \(e^{-1}\) with
\[
ee^{-1}=1.
\]
For every \(r\in R\),
\[
re=e.
\]
Multiplying by \(e^{-1}\),
\[
r=1.
\]
Thus all supported elements collapse to a single element \(1\). The unsupported element \(\tau\) remains the additive zero and multiplicative absorber. Therefore the localized semiring has exactly two elements:
\[
\tau,\qquad 1.
\]
The operations are precisely those of the Boolean semiring:
\[
1\oplus1=1,\qquad 1\oplus\tau=1,\qquad \tau\oplus\tau=\tau,
\]
and
\[
1\cdot1=1,\qquad 1\cdot\tau=\tau,\qquad \tau^2=\tau.
\]
\end{proof}

\section{Products and multiplicative reconstruction}

\begin{theorem}[Boolean fiber product]
For commutative rings \(A,B\), there is a canonical semiring isomorphism
\[
G(A\times B)\cong G(A)\times_{\BB}G(B),
\]
where the fiber product is taken over the support characters
\[
\chi_A:G(A)\to\BB,\qquad \chi_B:G(B)\to\BB.
\]
\end{theorem}

\begin{proof}
The fiber product consists of pairs \((x,y)\in G(A)\times G(B)\) such that
\[
\chi_A(x)=\chi_B(y).
\]
Thus either both are unsupported, giving \((\tau,\tau)\), or both are supported, giving a pair \((a,b)\in A\times B\).

Define
\[
\Phi:G(A\times B)\to G(A)\times_{\BB}G(B)
\]
by
\[
\Phi(\tau)=(\tau,\tau),
\]
and
\[
\Phi((a,b))=(a,b)
\]
with the right side read as the pair of supported elements \(a\in G(A)\), \(b\in G(B)\).

Define the inverse
\[
\Psi:G(A)\times_{\BB}G(B)\to G(A\times B)
\]
by
\[
\Psi(\tau,\tau)=\tau,
\]
and
\[
\Psi(a,b)=(a,b)\in A\times B
\]
on supported pairs.

Both maps preserve addition and multiplication componentwise, and they are inverse to one another.
\end{proof}

\begin{theorem}[Multiplicative reconstruction requires additive relations]
Let \(M_S\) be the multiplicative monoid of \(S=G(R)\). Let \(\NN[M_S]\) be the free commutative semiring on the monoid. Let \(\mathcal R\) be the semiring congruence generated by
\[
[a]+[b]\sim[a+b]\qquad(a,b\in R),
\]
and
\[
[x]+[\tau]\sim[x]\qquad(x\in S).
\]
Then
\[
\NN[M_S]/\mathcal R\cong S.
\]
\end{theorem}

\begin{proof}
The monoid map
\[
M_S\to S
\]
given by inclusion induces a semiring homomorphism
\[
\Theta:\NN[M_S]\to S.
\]
The defining relations of \(\mathcal R\) hold in \(S\), so \(\Theta\) descends to
\[
\overline\Theta:\NN[M_S]/\mathcal R\to S.
\]

The map \(\overline\Theta\) is surjective because \(S\) is generated by its underlying multiplicative monoid.

Every element of \(\NN[M_S]\) is a finite formal sum
\[
[x_1]+\cdots+[x_n].
\]
Using the relations \([x]+[\tau]\sim[x]\), all \([\tau]\)-summands are deleted unless every summand is \([\tau]\), in which case the sum is \([\tau]\).

If all remaining \(x_i\) are supported elements of \(R\), the relations
\[
[a]+[b]\sim[a+b]
\]
reduce the sum to a single symbol
\[
[x_1+\cdots+x_n]\in R.
\]
Thus every congruence class has a unique representative either \([\tau]\) or \([r]\) with \(r\in R\). These representatives multiply and add exactly as in \(G(R)\). Hence \(\overline\Theta\) is injective.
\end{proof}

\section{Semimodule classification}

Let \(S=G(R)\), and let \(M\) be a left \(S\)-semimodule.

\begin{definition}[Support projector]
Define
\[
\varepsilon_M:M\to M
\]
by
\[
\varepsilon_M(m):=em.
\]
Define the support skeleton
\[
L_M:=eM.
\]
For \(l\in L_M\), define the support fiber
\[
M_l:=\{m\in M:em=l\}.
\]
\end{definition}

\begin{lemma}[Support skeleton]
The set \(L_M\) is a join-semilattice with zero. The join is
\[
l\vee l':=l+l',
\]
and the zero is \(0_M\).
\end{lemma}

\begin{proof}
Let \(l=em\), \(l'=en\). Then
\[
l+l'=em+en=e(m+n),
\]
so \(l+l'\in eM=L_M\).

The operation is associative and commutative because addition in \(M\) is. It is idempotent because
\[
l+l=em+em=(e\oplus e)m=em=l,
\]
since
\[
e\oplus e=0_R+0_R=0_R=e.
\]

Finally,
\[
0_M=e0_M,
\]
so \(0_M\in L_M\), and \(0_M\) is the least element.
\end{proof}

\begin{lemma}[Fibers are modules]
Each fiber \(M_l\) is an \(R\)-module with zero element \(l\).
\end{lemma}

\begin{proof}
Let \(m\in M_l\). Then \(em=l\).

First,
\[
m+l=m+em=(1\oplus e)m=1m=m,
\]
because
\[
1\oplus e=1_R+0_R=1_R.
\]
Thus \(l\) is the additive zero inside \(M_l\).

If \(m,n\in M_l\), then
\[
e(m+n)=em+en=l+l=l,
\]
so \(m+n\in M_l\).

If \(r\in R\), then
\[
e(rm)=(er)m.
\]
But
\[
er=0_Rr=0_R=e,
\]
so
\[
e(rm)=em=l.
\]
Thus \(rm\in M_l\).

The additive inverse of \(m\in M_l\) is \((-1)m\), since
\[
m+(-1)m=(1\oplus(-1))m=e m=l.
\]
All module axioms follow from the semimodule axioms restricted to supported scalars \(R\subset S\).
\end{proof}

\begin{definition}[Transport maps]
If \(l\le l'\) in \(L_M\), meaning
\[
l\vee l'=l',
\]
define
\[
T_{l,l'}:M_l\to M_{l'}
\]
by
\[
T_{l,l'}(m)=m+l'.
\]
\end{definition}

\begin{lemma}[Transport maps are functorial and \(R\)-linear]
For \(l\le l'\le l''\),
\[
T_{l,l}=\Id_{M_l},
\]
and
\[
T_{l',l''}\circ T_{l,l'}=T_{l,l''}.
\]
Each \(T_{l,l'}\) is \(R\)-linear.
\end{lemma}

\begin{proof}
If \(m\in M_l\), then
\[
e(m+l')=em+el'=l+l'=l',
\]
so \(T_{l,l'}(m)\in M_{l'}\).

If \(l=l'\), then
\[
T_{l,l}(m)=m+l=m
\]
because \(l\) is the zero of \(M_l\).

For \(l\le l'\le l''\),
\[
T_{l',l''}(T_{l,l'}(m))
=
(m+l')+l''
=
m+(l'+l'')
=
m+l''
=
T_{l,l''}(m).
\]

For additivity, if \(m,n\in M_l\), then
\[
T_{l,l'}(m+n)=(m+n)+l'.
\]
On the other hand,
\[
T_{l,l'}(m)+T_{l,l'}(n)
=
(m+l')+(n+l')
=
m+n+l'+l'
=
m+n+l',
\]
because \(l'+l'=l'\).

For \(r\in R\),
\[
T_{l,l'}(rm)=rm+l'.
\]
Since \(l'=ey\) for some \(y\), one has
\[
rl'=r(ey)=(re)y=ey=l'.
\]
Therefore
\[
rT_{l,l'}(m)=r(m+l')=rm+rl'=rm+l'=T_{l,l'}(rm).
\]
\end{proof}

\begin{definition}[Support diagrams]
A support diagram over \(R\) is a triple
\[
A=(L,\{A_l\}_{l\in L},\{\rho_{l,l'}\}_{l\le l'})
\]
where:
\begin{enumerate}[label=(\roman*)]
\item \(L\) is a join-semilattice with zero;
\item each \(A_l\) is an \(R\)-module;
\item each \(\rho_{l,l'}:A_l\to A_{l'}\) is \(R\)-linear;
\item \(\rho_{l,l}=\Id\), and \(\rho_{l',l''}\rho_{l,l'}=\rho_{l,l''}\) whenever \(l\le l'\le l''\).
\end{enumerate}
Let \(\Diag_R\) denote the category of such diagrams.
\end{definition}

\begin{theorem}[From diagrams to semimodules]
Let
\[
A=(L,A_l,\rho_{l,l'})
\]
be a support diagram. Define
\[
M(A):=\bigsqcup_{l\in L}A_l.
\]
For \(x\in A_l\), \(y\in A_m\), define
\[
x\boxplus y
:=
\rho_{l,l\vee m}(x)+\rho_{m,l\vee m}(y)
\in A_{l\vee m}.
\]
For \(r\in R\), put
\[
r\cdot x:=rx\in A_l,
\]
and
\[
\tau\cdot x:=0_{A_0}.
\]
Then \(M(A)\) is a left \(G(R)\)-semimodule.
\end{theorem}

\begin{proof}
Commutativity of \(\boxplus\) follows from commutativity of addition in \(A_{l\vee m}\).

For associativity, take \(x\in A_l\), \(y\in A_m\), \(z\in A_n\), and put
\[
u=l\vee m\vee n.
\]
Using functoriality of the \(\rho\)'s,
\[
(x\boxplus y)\boxplus z
=
\rho_{l,u}(x)+\rho_{m,u}(y)+\rho_{n,u}(z),
\]
and the same expression is obtained from
\[
x\boxplus(y\boxplus z).
\]
Thus addition is associative.

The additive zero is \(0_{A_0}\). If \(x\in A_l\), then
\[
x\boxplus0_{A_0}
=
\rho_{l,l}(x)+\rho_{0,l}(0)=x.
\]

Scalar associativity and distributivity for supported \(r,s\in R\) follow fiberwise from the \(R\)-module structure and \(R\)-linearity of the transition maps.

The rules involving \(\tau\) also hold. Since \(\tau\) is the additive zero of \(G(R)\),
\[
(r\oplus\tau)\cdot x=r\cdot x.
\]
On the other side,
\[
r\cdot x\boxplus \tau\cdot x
=
r x\boxplus 0_{A_0}=rx.
\]
Similarly
\[
\tau\cdot(x\boxplus y)=0_{A_0}
\]
and
\[
\tau\cdot x\boxplus\tau\cdot y=0_{A_0}\boxplus0_{A_0}=0_{A_0}.
\]
Multiplicativity involving \(\tau\) follows from \(\tau r=\tau\) and \(\tau^2=\tau\).

Thus \(M(A)\) is a \(G(R)\)-semimodule.
\end{proof}

\begin{theorem}[Semimodule classification]
The assignments
\[
M\longmapsto
(L_M,\{M_l\}_{l\in L_M},\{T_{l,l'}\})
\]
and
\[
A\longmapsto M(A)
\]
define quasi-inverse equivalences
\[
G(R)\text{-}\mathbf{SMod}
\simeq
\Diag_R.
\]
\end{theorem}

\begin{proof}
Starting with a \(G(R)\)-semimodule \(M\), the preceding lemmas construct a support diagram.

Starting with a support diagram \(A\), the preceding theorem constructs a \(G(R)\)-semimodule.

We check the constructions are inverse.

Let \(M\) be a \(G(R)\)-semimodule. The fibers \(M_l\) partition \(M\). In the reconstructed semimodule, if \(m\in M_l\), \(n\in M_{l'}\), the reconstructed sum is
\[
T_{l,l\vee l'}(m)+T_{l',l\vee l'}(n)
=
(m+l\vee l')+(n+l\vee l').
\]
Since \(l\vee l'\) is the zero of the fiber containing \(m+n\), this equals
\[
m+n.
\]
Thus the reconstructed addition is the original addition. Scalar multiplication is similarly unchanged.

Conversely, start with a diagram \(A\). In \(M(A)\), the support of \(x\in A_l\) is \(l\), because the action of \(e=0_R\) sends \(x\) to the zero of the fiber \(A_l\), which is indexed by \(l\). The fiber over \(l\) is exactly \(A_l\), and the transport map is exactly \(\rho_{l,l'}\). Hence the recovered diagram is \(A\).

Functoriality is immediate: a \(G(R)\)-linear map preserves \(eM\), fibers, and transport maps; conversely, a morphism of diagrams preserves the reconstructed addition and scalar multiplication.
\end{proof}

\begin{corollary}[Cancellative sector]
The cancellative \(G(R)\)-semimodules are exactly ordinary \(R\)-modules under restriction along
\[
p_R:G(R)\to R.
\]
That is,
\[
G(R)\text{-}\mathbf{SMod}_{\mathrm{canc}}
\simeq
R\text{-}\mathbf{Mod}.
\]
\end{corollary}

\begin{proof}
Let \(M\) be cancellative. For any \(m\in M\),
\[
m+em=(1\oplus e)m=1m=m.
\]
Also
\[
m+0_M=m.
\]
By cancellativity,
\[
em=0_M.
\]
Thus
\[
eM=\{0_M\}.
\]
So the support diagram has one fiber, which is an \(R\)-module.

Conversely, any \(R\)-module becomes a \(G(R)\)-semimodule via
\[
r\cdot m=rm,\qquad \tau\cdot m=0.
\]
Its additive monoid is an abelian group, hence cancellative.
\end{proof}

\begin{corollary}[Free semimodule as Boolean cube]
The free semimodule \(G(R)^n\) has support diagram
\[
I\subseteq\{1,\dots,n\}
\quad\longmapsto\quad
R^I.
\]
Thus its support skeleton is the Boolean lattice
\[
\mathcal P(\{1,\dots,n\}).
\]
\end{corollary}

\begin{proof}
For
\[
x=(x_1,\dots,x_n)\in G(R)^n,
\]
the support element \(ex\) records exactly which coordinates are supported:
\[
ex_i=
\begin{cases}
\tau,&x_i=\tau,\\
e,&x_i\in R.
\end{cases}
\]
Thus supports are subsets of \(\{1,\dots,n\}\). The fiber over \(I\) consists of vectors supported exactly on \(I\), identified with \(R^I\). If \(I\subseteq J\), transport is extension by supported zeros in the new coordinates.
\end{proof}

\section{Projective semimodules}

\begin{definition}[Decorated finite-poset object]
Let \(J\) be a finite poset. Let \(Q_j\) be a finitely generated projective \(R\)-module for each \(j\in J\). Let \(\operatorname{Idl}(J)\) denote the lattice of order ideals of \(J\). Define a support diagram
\[
\mathcal P(J,Q)_U:=\bigoplus_{j\in U}Q_j
\]
for \(U\in\operatorname{Idl}(J)\), with transition maps given by direct-summand inclusions.
\end{definition}

\begin{theorem}[Finitely generated projective classification]
A finitely generated \(G(R)\)-semimodule is projective if and only if it is isomorphic to
\[
\mathcal P(J,Q)
\]
for a finite poset \(J\) and finitely generated projective \(R\)-modules \(Q_j\).

Moreover,
\[
\mathcal P(J,Q)\cong\mathcal P(J',Q')
\]
if and only if there is a poset isomorphism
\[
\phi:J\to J'
\]
and \(R\)-module isomorphisms
\[
Q_j\cong Q'_{\phi(j)}.
\]
\end{theorem}

\begin{proof}
First assume \(P=\mathcal P(J,Q)\). Since each \(Q_j\) is finitely generated projective, choose a finite free \(R\)-module \(R^{n_j}\) and a splitting
\[
Q_j\hookrightarrow R^{n_j}\twoheadrightarrow Q_j.
\]
The support lattice \(\operatorname{Idl}(J)\) is a retract of the Boolean lattice \(\mathcal P(J)\): include an ideal as a subset, and send a subset \(A\subseteq J\) to its downward closure \(\downarrow A\). These maps preserve finite joins.

Together with the module splittings, \(\mathcal P(J,Q)\) is a retract of a finite free Boolean-cube diagram, hence a retract of a finite free \(G(R)\)-semimodule. Therefore it is finitely generated projective.

Conversely, let \(P\) be finitely generated projective. Then \(P\) is a retract of \(G(R)^n\). Passing to support diagrams, \(eP\) is a finite join-retract of the Boolean lattice \(\mathcal P([n])\). Finite join-retracts of Boolean lattices are finite distributive lattices, hence by Birkhoff duality have the form
\[
\operatorname{Idl}(J)
\]
for a finite poset \(J\).

Each fiber \(P_U\) is a direct summand of a finite free \(R\)-module, hence finitely generated projective. For \(j\in J\), define the layer
\[
Q_j:=
P_{\downarrow j}/P_{(\downarrow j)\setminus\{j\}}.
\]
The transport maps split because \(P\) is a retract of a free diagram, so \(Q_j\) is finitely generated projective. Induction over order ideals gives
\[
P_U\cong\bigoplus_{j\in U}Q_j,
\]
with transition maps identified with direct-summand inclusions. Thus
\[
P\cong\mathcal P(J,Q).
\]

Finally, the support skeleton is intrinsic. An isomorphism of semimodules induces an isomorphism of support lattices
\[
\operatorname{Idl}(J)\cong\operatorname{Idl}(J'),
\]
hence a poset isomorphism \(J\cong J'\). The quotient layer
\[
P_{\downarrow j}/P_{(\downarrow j)\setminus\{j\}}
\]
is intrinsic, so the corresponding \(Q_j\)'s are isomorphic. The converse is immediate.
\end{proof}

\section{Conservative ideal and zeta side}

\begin{theorem}[Ideal multiplication is classical]
For ideals \(J,K\triangleleft R\),
\[
I_JI_K=I_{JK},
\]
where
\[
I_J=J\cup\{\tau\}.
\]
\end{theorem}

\begin{proof}
Every product of an element of \(I_J\) and an element of \(I_K\) is either \(\tau\), or a supported product \(jk\) with \(j\in J\), \(k\in K\). Therefore
\[
I_JI_K\subseteq I_{JK}.
\]

Conversely, \(\tau\in I_JI_K\) because \(\tau=\tau\cdot\tau\). Every element of \(JK\) is a finite sum of products \(j_ak_a\), with \(j_a\in J\), \(k_a\in K\). Each product lies in \(I_JI_K\), and the ideal is closed under addition. Hence
\[
JK\subseteq I_JI_K.
\]
Thus
\[
I_JI_K=I_{JK}.
\]
\end{proof}

\begin{corollary}[Class group preservation]
If \(R=\mathcal O_K\) is the ring of integers of a number field, then the ideal class monoid of nontrivial ideals of \(G(\mathcal O_K)\) is canonically the ordinary ideal class monoid of \(\mathcal O_K\). In particular,
\[
\Cl(G(\mathcal O_K))\cong\Cl(\mathcal O_K).
\]
\end{corollary}

\begin{proof}
The map
\[
J\longmapsto J\cup\{\tau\}
\]
is a multiplicative bijection from ordinary nonzero ideals of \(\mathcal O_K\) to nontrivial finite-norm ideals of \(G(\mathcal O_K)\). By the theorem, it preserves multiplication. Principal ideals are preserved because
\[
(a)\mapsto(a)\cup\{\tau\}.
\]
Therefore the class monoids are canonically isomorphic.
\end{proof}

\begin{corollary}[Dedekind zeta preservation]
For a number field \(K\),
\[
\zeta_{G(\mathcal O_K)}(s)=\zeta_K(s)
\]
on the normed ideal sheet.
\end{corollary}

\begin{proof}
The nontrivial finite-norm ideals of \(G(\mathcal O_K)\) are \(I_J=J\cup\{\tau\}\), with \(J\) a nonzero ideal of \(\mathcal O_K\). The quotient norm is the ordinary ideal norm
\[
N(I_J)=N(J).
\]
Therefore
\[
\zeta_{G(\mathcal O_K)}(s)
=
\sum_{0\neq J\triangleleft\mathcal O_K}N(J)^{-s}
=
\zeta_K(s).
\]
\end{proof}

\begin{theorem}[Separated quotient-size weight is nonmultiplicative]
Assume \(R\) has a multiplicative ideal norm \(N\). Define
\[
w(J):=N(J)+1.
\]
Then
\[
w(JK)=N(J)N(K)+1,
\]
but
\[
w(J)w(K)=N(J)N(K)+N(J)+N(K)+1.
\]
Thus \(w\) is not multiplicative for nonzero finite-norm ideals.
\end{theorem}

\begin{proof}
Since \(N\) is multiplicative,
\[
N(JK)=N(J)N(K).
\]
Therefore
\[
w(JK)=N(JK)+1=N(J)N(K)+1.
\]
But
\[
w(J)w(K)
=
(N(J)+1)(N(K)+1)
\]
\[
=
N(J)N(K)+N(J)+N(K)+1.
\]
Since \(N(J),N(K)>0\), these differ. Hence \(w\) is not multiplicative.
\end{proof}

\begin{corollary}[Separated zeta sheet is non-Eulerian]
The Dirichlet series
\[
\zeta_{G(R)}^{\mathrm{sep}}(s)
:=
\sum_{0\neq J}(N(J)+1)^{-s}
\]
has no Euler product indexed by the ordinary prime ideals.
\end{corollary}

\begin{proof}
An Euler product indexed by prime ideals requires a multiplicative coefficient system on ideals. The coefficient weight \(J\mapsto N(J)+1\) is not multiplicative by the preceding theorem. Hence no ordinary Euler product exists for this separated sheet.
\end{proof}

\section{Shifted zeta deformation}

Let \(L(s)\) be a Dirichlet series in a right half-plane. Define the shift operator
\[
(Tf)(s)=f(s+1),
\]
and the operator
\[
D=sT.
\]

\begin{definition}[Shifted flow]
Define formally
\[
L_t:=e^{-tD}L.
\]
\end{definition}

\begin{theorem}[Universal shifted-flow jets]
For all \(n\ge0\),
\[
\partial_t^n L_t(s)
=
(-1)^n(s)_nL_t(s+n),
\]
where
\[
(s)_n=s(s+1)\cdots(s+n-1).
\]
\end{theorem}

\begin{proof}
Since
\[
L_t=e^{-tD}L,
\]
one has
\[
\partial_t^nL_t=(-D)^ne^{-tD}L=(-1)^nD^nL_t.
\]
It remains to compute \(D^n\). Because
\[
D=sT,
\]
we have
\[
D^2f(s)=sT(sf(s+1))=s(s+1)f(s+2).
\]
Inductively,
\[
D^nf(s)=(s)_nf(s+n).
\]
Therefore
\[
\partial_t^nL_t(s)=(-1)^n(s)_nL_t(s+n).
\]
\end{proof}

\begin{theorem}[Hurwitz specialization]
For
\[
L(s)=\zeta(s),
\]
the shifted flow gives
\[
L_t(s)=\zeta(s,1+t).
\]
In particular, for integer \(q\ge0\),
\[
L_q(s)=\zeta(s)-\sum_{n=1}^{q}n^{-s}.
\]
\end{theorem}

\begin{proof}
The Hurwitz zeta function satisfies
\[
\partial_t\zeta(s,1+t)
=
-s\zeta(s+1,1+t).
\]
This is exactly the differential equation
\[
\partial_tL_t=-DL_t.
\]
At \(t=0\),
\[
\zeta(s,1)=\zeta(s).
\]
Thus uniqueness of the formal flow gives
\[
L_t(s)=\zeta(s,1+t).
\]

For integer \(q\ge0\),
\[
\zeta(s,1+q)
=
\sum_{n=0}^\infty(n+1+q)^{-s}
=
\sum_{m=q+1}^\infty m^{-s}
=
\zeta(s)-\sum_{m=1}^q m^{-s}.
\]
\end{proof}

\begin{theorem}[Endpoint non-Eulerianity]
For \(0<t<1\), the shifted Hurwitz sheet
\[
\zeta(s,1+t)
\]
is not an ordinary Dirichlet series in integer powers \(n^{-s}\). At \(t=1\),
\[
\zeta(s,2)=\zeta(s)-1,
\]
which is an ordinary Dirichlet series but has no ordinary Euler product.
\end{theorem}

\begin{proof}
Let \(\sigma\to+\infty\). The leading exponential scale of
\[
\zeta(\sigma,1+t)
\]
is
\[
(1+t)^{-\sigma}.
\]
If it were an ordinary Dirichlet series
\[
\sum_{n\ge1}a_n n^{-s},
\]
then the leading scale would be \(n_0^{-\sigma}\), where \(n_0\in\NN\) is the least \(n\) with \(a_n\neq0\). Thus
\[
n_0=1+t.
\]
For \(0<t<1\), this is impossible because \(1+t\notin\NN\).

At \(t=1\),
\[
\zeta(s,2)=\sum_{n\ge2}n^{-s}=\zeta(s)-1.
\]
A standard Euler product has coefficient \(1\) at \(1^{-s}\), coming from the empty product. But \(\zeta(s)-1\) has coefficient \(0\) at \(1^{-s}\). Hence it has no ordinary Euler product.
\end{proof}

\section{Boundary coprimality}

The split-zero correction forces the supported arithmetic zero \(e=0\) and the unit \(1\) to be treated as boundary states, not suppressed as degenerate exceptions.

\begin{definition}[Boundary coprimality matrices]
For \(q\ge1\), define
\[
B_q^+:=\{1,2,\dots,q\}.
\]
Let
\[
C_q^+=(c_{ij})_{1\le i,j\le q}
\]
where
\[
c_{ij}=
\begin{cases}
1,&\gcd(i,j)=1,\\
0,&\gcd(i,j)>1.
\end{cases}
\]

Define
\[
B_q^0:=\{0,1,\dots,q\}.
\]
Let
\[
C_q^0=(c_{ij})_{0\le i,j\le q}
\]
with the classical conventions
\[
\gcd(0,0)=0,\qquad \gcd(0,n)=n.
\]
Thus
\[
0\perp n
\quad\Longleftrightarrow\quad
n=1.
\]
\end{definition}

\begin{theorem}[Nondegenerate \(0/1/2/3\) boundary]
The positive matrix \(C_q^+\) is unimodular exactly for
\[
q=1,2,3.
\]
For \(q=3\),
\[
C_3^+
=
\begin{pmatrix}
1&1&1\\
1&0&1\\
1&1&0
\end{pmatrix},
\qquad
(C_3^+)^{-1}
=
\begin{pmatrix}
-1&1&1\\
1&-1&0\\
1&0&-1
\end{pmatrix}.
\]

The supported-zero matrix \(C_q^0\) is nontrivially unimodular precisely at \(q=3\). For \(q=3\),
\[
C_3^0=
\begin{pmatrix}
0&1&0&0\\
1&1&1&1\\
0&1&0&1\\
0&1&1&0
\end{pmatrix},
\]
with inverse
\[
(C_3^0)^{-1}=
\begin{pmatrix}
1&1&-1&-1\\
1&0&0&0\\
-1&0&0&1\\
-1&0&1&0
\end{pmatrix}.
\]
\end{theorem}

\begin{proof}
For \(q=1\),
\[
C_1^+=(1),
\]
so the determinant is \(1\).

For \(q=2\),
\[
C_2^+=
\begin{pmatrix}
1&1\\
1&0
\end{pmatrix},
\]
whose determinant is \(-1\).

For \(q=3\), direct computation gives
\[
\det
\begin{pmatrix}
1&1&1\\
1&0&1\\
1&1&0
\end{pmatrix}
=1,
\]
and multiplication verifies the displayed inverse.

For \(q\ge4\), the rows indexed by \(2\) and \(4\) in \(C_q^+\) coincide, because
\[
\gcd(n,2)=1
\quad\Longleftrightarrow\quad
n\text{ is odd}
\]
and
\[
\gcd(n,4)=1
\quad\Longleftrightarrow\quad
n\text{ is odd}.
\]
Thus \(C_q^+\) is singular for \(q\ge4\).

For \(C_q^0\), if \(q=2\), the rows indexed by \(0\) and \(2\) both have support only at the unit column \(1\), so the matrix is singular. For \(q\ge4\), the rows indexed by \(2\) and \(4\) are again identical. For \(q=3\), direct multiplication verifies the displayed inverse. Thus \(q=3\) is the unique nontrivial supported-zero unimodular boundary.
\end{proof}

\section{Residual divisor certificates: split-zero repair}

This section gives the corrected split-zero version of the residual divisor coefficient used in Erdős--Straus shell analysis.

Let \(p\) be prime. Fix a shell
\[
a=a_R(p)=\frac{p+R}{4},
\qquad
N=pa.
\]
Let
\[
g=\gcd(R,N),\qquad R_0=\frac Rg,\qquad N_0=\frac Ng.
\]
Then
\[
\gcd(R_0,N_0)=1.
\]

\begin{theorem}[Reduced divisor criterion]
The shell \(R\) gives a completion
\[
\frac4p=\frac1a+\frac1b+\frac1c
\]
if and only if there exists a divisor
\[
d\mid N_0^2
\]
such that
\[
d\equiv -N_0\pmod{R_0}.
\]
Equivalently,
\[
dN_0^{-1}\equiv -1\pmod{R_0}.
\]
\end{theorem}

\begin{proof}
The residual equation is
\[
\frac4p-\frac1a=\frac{4a-p}{pa}=\frac RN=\frac{R_0}{N_0}.
\]
Thus the completion condition is
\[
\frac{R_0}{N_0}=\frac1b+\frac1c.
\]
Multiplying by \(N_0bc\) gives
\[
R_0bc=N_0(b+c).
\]
Equivalently,
\[
(R_0b-N_0)(R_0c-N_0)=N_0^2.
\]
Set
\[
d=R_0b-N_0,\qquad e=R_0c-N_0.
\]
Then
\[
de=N_0^2,
\]
and
\[
d\equiv -N_0\pmod{R_0}.
\]

Conversely, suppose \(d\mid N_0^2\) and
\[
d\equiv -N_0\pmod{R_0}.
\]
Put
\[
e=\frac{N_0^2}{d}.
\]
Since \(\gcd(R_0,N_0)=1\), divisors of \(N_0^2\) are units modulo \(R_0\). Thus
\[
e\equiv \frac{N_0^2}{-N_0}\equiv -N_0\pmod{R_0}.
\]
Therefore
\[
b=\frac{N_0+d}{R_0},
\qquad
c=\frac{N_0+e}{R_0}
\]
are positive integers. The identity
\[
de=N_0^2
\]
then expands back to
\[
R_0bc=N_0(b+c),
\]
so
\[
\frac{R_0}{N_0}=\frac1b+\frac1c.
\]
The normalized congruence follows by multiplying by \(N_0^{-1}\) modulo \(R_0\).
\end{proof}

Factor
\[
N_0=\prod_{j=1}^s q_j^{e_j}.
\]
Define the signed exponent box
\[
B(N_0):=\prod_{j=1}^s[-e_j,e_j]_{\ZZ}.
\]
Define
\[
\lambda:B(N_0)\to(\ZZ/R_0\ZZ)^\times
\]
by
\[
\lambda(\beta_1,\dots,\beta_s)=\prod_{j=1}^s q_j^{\beta_j}.
\]

\begin{theorem}[Signed exponent criterion]
The shell works if and only if
\[
-1\in\lambda(B(N_0)).
\]
\end{theorem}

\begin{proof}
Every divisor \(d\mid N_0^2\) is uniquely
\[
d=\prod_{j=1}^s q_j^{k_j},
\qquad
0\le k_j\le 2e_j.
\]
Then
\[
dN_0^{-1}
=
\prod_{j=1}^s q_j^{k_j-e_j}.
\]
Set
\[
\beta_j=k_j-e_j.
\]
Then
\[
-e_j\le \beta_j\le e_j,
\]
so \(\beta\in B(N_0)\), and
\[
dN_0^{-1}=\lambda(\beta).
\]
By the reduced divisor criterion, the shell works exactly when some divisor has
\[
dN_0^{-1}\equiv -1.
\]
This is equivalent to
\[
-1\in\lambda(B(N_0)).
\]
\end{proof}

\begin{definition}[Split-zero residual coefficient]
Let
\[
K:=\langle q_1,\dots,q_s\rangle\le(\ZZ/R_0\ZZ)^\times.
\]
Let \(G(\NN)\) be the split-zero semiring built from \(\NN\). Define
\[
\mathfrak C_{R,p}(t)\in G(\NN)
\]
by
\[
\mathfrak C_{R,p}(t)
=
\begin{cases}
\#\{\beta\in B(N_0):\lambda(\beta)=t\},&
\text{if this number is positive},\\
\tau,&
\text{if this number is zero}.
\end{cases}
\]
\end{definition}

\begin{theorem}[Support-complete shell criterion]
The shell \(R\) works if and only if
\[
\mathfrak C_{R,p}(-1)\neq\tau.
\]
Equivalently,
\[
\chi_{\NN}(\mathfrak C_{R,p}(-1))=1.
\]
The ordinary integer count is the collapsed shadow
\[
p_{\NN}(\mathfrak C_{R,p}(-1)).
\]
\end{theorem}

\begin{proof}
By definition,
\[
\mathfrak C_{R,p}(-1)\neq\tau
\]
if and only if
\[
\#\{\beta\in B(N_0):\lambda(\beta)=-1\}>0.
\]
This holds if and only if
\[
-1\in\lambda(B(N_0)).
\]
By the signed exponent criterion, this is exactly the condition that the shell works.

The equality for the ordinary count follows because the ring reflection
\[
p_{\NN}:G(\NN)\to\NN
\]
sends \(\tau\mapsto0\) and supported counts to themselves.
\end{proof}

\section{Formal repair rules}

The previous theorems imply the following mandatory repair protocol.

\begin{theorem}[Zero repair]
Every assertion over \(G(R)\) containing
\[
x=0
\]
must be replaced by one of:
\[
x=\tau,
\]
\[
x=e=0_R,
\]
or
\[
p_R(x)=0_R.
\]
These are pairwise non-equivalent conditions in general.
\end{theorem}

\begin{proof}
The fiber of \(p_R\) over \(0_R\) is
\[
\{\tau,e\}.
\]
The support character distinguishes them:
\[
\chi_R(\tau)=0,
\qquad
\chi_R(e)=1.
\]
Therefore \(p_R(x)=0_R\), \(x=\tau\), and \(x=e\) are distinct conditions.
\end{proof}

\begin{theorem}[Nonzero repair]
Every assertion over \(G(R)\) containing
\[
x\neq0
\]
must specify whether the intended condition is:
\[
x\neq\tau,
\]
\[
\chi_R(x)=1,
\]
\[
x\in R,
\]
or
\[
x\in R\setminus\{0_R\}.
\]
Only the last condition is ordinary arithmetic nonzero-ness.
\end{theorem}

\begin{proof}
One has
\[
x\neq\tau
\quad\Longleftrightarrow\quad
x\in R
\quad\Longleftrightarrow\quad
\chi_R(x)=1.
\]
But this allows
\[
x=e=0_R.
\]
Thus ordinary arithmetic nonzero-ness is the strictly stronger condition
\[
x\in R\setminus\{0_R\}.
\]
\end{proof}

\begin{theorem}[Prime repair]
Every statement using the word ``prime'' over \(G(R)\) must specify which of the following is meant:
\begin{enumerate}[label=(\roman*)]
\item the support prime ideal \(P_\tau=\{\tau\}\);
\item an ordinary supported prime ideal \(P_{\mathfrak p}=\mathfrak p\cup\{\tau\}\);
\item the supported-zero prime ideal \(P_0=\{0_R,\tau\}\) when \(R\) is a domain;
\item a nonzero supported irreducible element of \(R\);
\item the divisibility-prime supported zero \(e=0_R\).
\end{enumerate}
\end{theorem}

\begin{proof}
The ideal alternatives follow from the prime ideal classification. The nonzero supported irreducible alternative follows from the nonzero irreducible theorem. The element \(e=0_R\) is divisibility-prime but reducible by the supported-zero prime theorem. Hence the word ``prime'' is ambiguous unless its support layer is specified.
\end{proof}

\begin{theorem}[Euler repair]
The classical ideal zeta sheet remains Eulerian:
\[
\zeta_{G(\mathcal O_K)}(s)=\zeta_K(s).
\]
The separated support quotient-size sheet is non-Eulerian:
\[
\sum_{0\neq J}(N(J)+1)^{-s}
\]
has no ordinary Euler product indexed by prime ideals.
\end{theorem}

\begin{proof}
The first statement is the Dedekind zeta preservation theorem. The second is the nonmultiplicativity of the separated weight \(N(J)+1\).
\end{proof}

\begin{theorem}[Linear algebra repair]
Every \(G(R)\)-linear proof must be carried out in the support-diagram category
\[
G(R)\text{-}\mathbf{SMod}
\simeq
\Diag_R.
\]
It may be collapsed to ordinary \(R\)-linear algebra only if the semimodule is cancellative.
\end{theorem}

\begin{proof}
This is exactly the semimodule classification theorem and its cancellative-sector corollary.
\end{proof}

\section{Conclusion}

The split-zero theory is conservative on the ordinary nonzero supported arithmetic sheet, but not at the boundary
\[
\{\tau,e,1_R\}.
\]
The exact repaired conclusions are:

\[
\boxed{
\text{ordinary nonzero supported factorization survives unchanged;}
}
\]

\[
\boxed{
e=0_R\text{ is reducible and divisibility-prime;}
}
\]

\[
\boxed{
\text{prime}\not\Rightarrow\text{irreducible in }G(R);
}
\]

\[
\boxed{
\text{ordinary ideal multiplication, ideal class monoids, and Dedekind zeta functions are preserved;}
}
\]

\[
\boxed{
\text{separated support quotient-size zeta sheets are non-Eulerian;}
}
\]

\[
\boxed{
G(R)\text{-linear algebra is support-diagrammatic, not ordinary module theory;}
}
\]

\[
\boxed{
\text{residual divisor coefficients must be lifted to split-zero coefficients before applying support logic.}
}
\]

Consequently, every proof formerly written over \(G(R)\) by collapsing through
\[
p_R:G(R)\to R
\]
proves only the classical ring-shadow theorem. A support-complete theorem must carry \(\tau\), \(e=0_R\), units, nonzero supported elements, support skeletons, and finite transport coefficients explicitly.

\begin{thebibliography}{9}

\bibitem{globalization}
The Clankers.
\emph{A Note on the Globalization Semiring \(G(\mathbb Z)\)}.
March 2026.

\bibitem{secondary}
The Clankers.
\emph{A Secondary Note on Split-Zero Globalization: Semimodule Classification, Multiplicative Reconstruction, and All-Orders Zeta Deformation}.
April 2026.

\bibitem{cue}
The Clankers.
\emph{Split-Zero CUE/Hurwitz Phase-Space Revised}.
April 2026.

\bibitem{support}
The Clankers.
\emph{Support-Idempotent Homogenization for a Split-Zero Collatz F-Series Frame}.
2026.

\bibitem{residual}
The Clankers.
\emph{Normalized Divisor Supports and Residual Divisor Certificates for the Erdős--Straus Equation}.
2026.

\end{thebibliography}

\end{document}